% ЗГЕНЕРОВАНО scripts/build_format_v2_chapter.py — не редагувати вручну.
% Джерело: editorial/cards/10.*.json, format-v2/rozdil-10.json, paraleli.json.

\naukchapter{10}{8}{Глава 10. Досить уже}
\naukintro{Десята глава — найрізкіша за тоном. Вона не пояснює й не доводить: вона перелічує те, що вже було, і каже досить. Спершу відкидаються бажання й користь (10.1), далі побутові речі названо сном і чарами (10.2), далі спрага названа самою сансарою (10.3–10.4). Останні три вірші звернені до минулого: усе це вже гинуло не раз, і саме тому питання стоїть не про нове, а про припинення.}
\vsec{10.1}
\vdeva{विहाय वैरिणं कामम् अर्थं चानर्थसंकुलं । धर्मम् अप्य् एतयोर् हेतुं सर्वत्रानादरं कुरु ॥}
\vtranslit{віхая вайрінам камам артхам чанартхасамкулам | дгармам апй етайор хетум сарватранадарам куру ||}
\vperek{Покинувши вороже бажання і сповнену лиха користь — навіть дхарму, що є причиною обох, — усюди зневажай.}
\vcom{Построкове зіставлення: редакторська чернетка, не рецензовано.}
\vcom{Рядок а}
\vdeva{विहाय वैरिणं कामम् अर्थं चानर्थसंकुलं}
\vtranslit{віхая вайрінам камам артхам чанартхасамкулам}
\vperek{Покинувши вороже бажання і сповнену лиха користь —}
\vcom{vihāya — дієприслівникова форма «покинувши»; vairiṇaṃ kāmam і arthaṃ cānarthasaṃkulaṃ називають бажання та користь, сповнену лиха. Цей епітет стосується artha, а не автоматично всіх наступних понять.}
\vcom{Рядок б}
\vdeva{धर्मम् अप्य् एतयोर् हेतुं सर्वत्रानादरं कुरु}
\vtranslit{дгармам апь етайор хетум сарватранадарам куру}
\vperek{навіть дхарму, що є причиною обох, — усюди зневажай.}
\vcom{dharmam api вводить навіть дхарму; etayor hetuṃ відносить «причину обох» до попередніх понять, sarvatrānādaraṃ kuru має наказову силу. Порівняння з українською етикою потребує контексту: це теза вірша, а не заклик до порушення закону.}
\vcom{Дхарма стоїть у переліку **останньою** і з підсиленням api («навіть»): спершу відкидають бажання й користь, а потім те, що їх породжує. anartha-saṃkula («сповнене лиха») стосується саме artha, і переклад не повинен переносити цю оцінку на дхарму. «Зневажай» (anādara) — сильне слово; м'якше «не зважай» змінило б ступінь.}
\vpar{Українська паралель (наша інтерпретація)}{Пантелеймон Куліш, «Хуторна поезія» (1882).}{Наша паралель: у Куліша відмова від користі теж подана як умова спокою, а не як втрата. Наше порівняння — типологічне, не твердження про джерело.}
\vfig{ilustratsii/10-1-neti.pdf}{Бажання, користь і дхарма відкинуті разом}{Схема наша, не історична ілюстрація.}

\vsec{10.2}
\vdeva{स्वप्नेन्द्रजालवत् पश्य दिनानि त्रीणि पञ्च वा । मित्रक्षेत्रधनागारदारदायादिसंपदः ॥}
\vtranslit{свапнендраджалават пашя дінані тріні паньча ва | мітракшетрадганаґарадарадаядісампадах ||}
\vperek{Дивись на друзів, поле, гроші, дім, дружину, спадок і решту багатств як на сон і чари — на три дні або п'ять.}
\vcom{Построкове зіставлення: редакторська чернетка, не рецензовано.}
\vcom{Рядок а}
\vdeva{स्वप्नेन्द्रजालवत् पश्य दिनानि त्रीणि पञ्च वा}
\vtranslit{свапнендраджалават пашья дінані тріні панча ва}
\vperek{Дивись на друзів, поле, гроші, дім, дружину, спадок і решту багатств}
\vcom{svapnendrajālavat paśya пропонує погляд «як на сон і чари», dināni trīṇi pañca vā згадує три чи п'ять днів. Український повний переклад переносить перелік предметів із наступного піврядка перед порівнянням; рядковий фрагмент тому є смисловим, а не буквально послідовним підрядником.}
\vcom{Рядок б}
\vdeva{मित्रक्षेत्रधनागारदारदायादिसंपदः}
\vtranslit{мітракшетрадганаґарадарадаядісампадах}
\vperek{як на сон і чари — на три дні або п'ять.}
\vcom{mitrakṣetradhanāgāradāradāyādisaṃpadaḥ складає конкретний перелік дружби, поля, грошей, житла, родини та спадку. В українській цілісній редакції цей перелік стоїть на початку речення; це прозоре редакторське переставлення, а не текстова розбіжність санскритських свідків.}
\vcom{Перелік навмисно побутовий: не «світ» загалом, а конкретні речі, за які тримаються. «Три дні або п'ять» — не точний строк, а спосіб сказати «коротко»; переклад зберігає числа, бо вони й творять ефект. indrajāla — не просто «ілюзія», а фокус, свідома підміна видимого.}
\vpar{Українська паралель (ПОТРЕБУЄ ПЕРЕВІРКИ)}{українська барокова проповідь, не визначено.}{ПОТРЕБУЄ ПЕРЕВІРКИ: перелік побутових речей як сну й чару близький до барокової проповіді про марність, але конкретного твору ми не перевірили.}
\vfig{ilustratsii/10-2-illusion.pdf}{Побутові речі під серпанком: сон і чари}{Схема наша, не історична ілюстрація.}

\vsec{10.3}
\vdeva{यत्र यत्र भवेत् तृष्णा संसारं विद्धि तत्र वै । प्रौढवैराग्यम् आश्रित्य वीततृष्णः सुखी भव ॥}
\vtranslit{ятра ятра бгавет трішна самсарам віддгі татра вай | праудгавайраґям ашрітя вітатрішнах сукхі бгава ||}
\vperek{Де б не виникла спрага — знай: саме там і сансара. Опираючись на зрілу безсторонність, стань тим, у кого спрага минула, — і будь щасливий.}
\vcom{Построкове зіставлення: редакторська чернетка, не рецензовано.}
\vcom{Рядок а}
\vdeva{यत्र यत्र भवेत् तृष्णा संसारं विद्धि तत्र वै}
\vtranslit{ятра ятра бгавет трішна самсарам віддгі татра вай}
\vperek{Де б не виникла спрага — знай: саме там і сансара.}
\vcom{Пара yatra yatra ... tatra vai утворює відповідність «де б ... саме там»; bhavet tṛṣṇā і saṃsāraṃ viddhi поєднують спрагу та сансару через наказ «знай». Тут сансара названа у зв'язку зі спрагою, а не як географічне місце.}
\vcom{Рядок б}
\vdeva{प्रौढवैराग्यम् आश्रित्य वीततृष्णः सुखी भव}
\vtranslit{праудгавайраґьям ашрітья вітатрішнах сукхі бгава}
\vperek{Опираючись на зрілу безсторонність, стань тим, у кого спрага минула, — і будь щасливий.}
\vcom{pra uḍha-vairāgyam āśritya описує опертя на зрілу неприв'язаність; vītatṛṣṇaḥ sukhī bhava завершується наказом «будь щасливий». Складене prauḍhavairāgyam збережено в IAST картки; коментар не додає нового варіанта читання.}
\vcom{Тут уперше в розділі дано не заперечення, а **вказівку**: `yatra yatra … tatra` («де б … саме там»). Сансара — не місце й не стан світу, а точка, де виникла спрага. prauḍha («зріла») важливе: безсторонність має бути витриманою, а не поривчастою. Висновок «будь щасливий» — наказ, а не побажання.}
\vpar{Українська паралель (наша інтерпретація)}{Володимир Винниченко, «Сонячна машина» (1928).}{Наша паралель: у Винниченка спрага теж постає тим, що тримає людину в колі, незалежно від того, чи вона цього хоче. Наше порівняння — типологічне.}
\vfig{ilustratsii/10-3-knot.pdf}{Вузол спраги: де вона виникла, там і сансара}{Схема наша, не історична ілюстрація.}

\vsec{10.4}
\vdeva{तृष्णामात्रात्मको बन्धस् तन्नाशो मोक्ष उच्यते । भवासंसक्तिमात्रेण प्राप्तितुष्टिर् मुहुर् मुहुः ॥}
\vtranslit{трішнаматратмако бандгас таннашо мокша учяте | бгавасамсактіматрена праптітуштір мухур мухух ||}
\vperek{Полон має природу самої лише спраги; її знищення називають визволенням. Від самої лише неприв'язаності до буття приходить вдоволеність досягнутим — знову й знову.}
\vcom{Построкове зіставлення: редакторська чернетка, не рецензовано.}
\vcom{Рядок а}
\vdeva{तृष्णामात्रात्मको बन्धस् तन्नाशो मोक्ष उच्यते}
\vtranslit{трішнаматратмако бандгас таннашо мокша учьяте}
\vperek{Полон має природу самої лише спраги; її знищення називають визволенням.}
\vcom{tṛṣṇāmātrātmako bandhas визначає полон через саму спрагу; tannāśo mokṣa ucyate називає її зникнення визволенням. Це визначення всередині вірша, не універсальна психологічна діагностика.}
\vcom{Рядок б}
\vdeva{भवासंसक्तिमात्रेण प्राप्तितुष्टिर् मुहुर् मुहुः}
\vtranslit{бгавасамсактіматрена праптітуштір мухур мухух}
\vperek{Від самої лише неприв'язаності до буття приходить вдоволеність досягнутим — знову й знову.}
\vcom{bhavāsaṃsaktimātreṇa указує на неприв'язаність, prāptituṣṭir — на вдоволеність досягнутим, muhur muhuḥ — на повторюваність. Українське «знову й знову» зберігає повтор без додавання нового джерела.}
\vcom{Тут полон означено **однією** ознакою — спрага, і визволення теж однією — її знищення. Це найточніше місце глави, і воно коротке; розширювати переклад не варто. `toṣa` входить сюди як складник `prāpti-tuṣṭi` («вдоволеність досягнутим»), а не як самостійна чеснота — це варто пам'ятати при звірці з 1.2, де `toṣa` стоїть окремо. `muhur muhuḥ` — не «завжди», а «раз у раз»: вдоволеність повторюється, а не тримається одного разу назавжди.}
\vpar{Українська паралель (наша інтерпретація)}{Іван Франко, «Смерть Каїна» (1889).}{Наша паралель: у Франка полон теж має природу самої лише спраги, а не зовнішньої перешкоди. Наше порівняння — за будовою думки.}
\vfig{ilustratsii/10-4-knot.pdf}{Один вузол у двох назвах: полон і визволення}{Схема наша, не історична ілюстрація.}

\vsec{10.5}
\vdeva{त्वम् एकश् चेतनः शुद्धो जडं विश्वम् असत् तथा । अविद्यापि न किञ्चित् सा का बुभुत्सा तथापि ते ॥}
\vtranslit{твам екаш четанах шуддго джадам вішвам асат татха | авідяпі на кіньчіт са ка бубгутса татхапі те ||}
\vperek{Ти один — свідомий, чистий; світ же — неживий і несправжній. Навіть незнання — аніщо. То яке в тебе, при всьому тому, бажання пізнавати?}
\vcom{Построкове зіставлення: редакторська чернетка, не рецензовано.}
\vcom{Рядок а}
\vdeva{त्वम् एकश् चेतनः शुद्धो जडं विश्वम् असत् तथा}
\vtranslit{твам екаш четанах шуддго джадам вішвам асат татха}
\vperek{Ти один — свідомий, чистий; світ же — неживий і несправжній.}
\vcom{tvam ekaś cetanaḥ śuddho характеризує адресата як свідомого й чистого, jaḍaṃ viśvam asat tathā — світ як неживий та несправжній у межах позиції твору. asat не дорівнює моральній категорії «злий».}
\vcom{Рядок б}
\vdeva{अविद्यापि न किञ्चित् सा का बुभुत्सा तथापि ते}
\vtranslit{авідьяпі на кінчіт са ка бубгутса татхапі те}
\vperek{Навіть незнання — аніщо. То яке в тебе, при всьому тому, бажання пізнавати?}
\vcom{avidyāpi na kiñcit sā — твердження про незнання, kā bubhutsā tathāpi te — риторичне питання про бажання пізнавати. Це філософська постановка питання, не заклик відмовлятися від навчання чи наукової перевірки.}
\vcom{`asat` тут — «несправжнє», не «зло»; це різні речі, і переклад не повинен їх зливати. Питання `kā bubhutsā` замикає коло: якщо навіть незнання ніщо, то пізнавати нічого — і саме тому питання адресоване тому, хто вже слухає пояснення. Це найгостріше місце глави: воно ставить під сумнів користь самої розмови.}
\vpar{Українська паралель (ПОТРЕБУЄ ПЕРЕВІРКИ)}{українська релігійна філософія, не визначено.}{ПОТРЕБУЄ ПЕРЕВІРКИ: питання, навіщо пізнавати те, що й так не є річчю, близьке до української релігійної філософії, але конкретного місця ми не звірили.}
\vfig{ilustratsii/10-5-witness.pdf}{Свідомий у центрі, несправжній світ навколо}{Схема наша, не історична ілюстрація.}

\vsec{10.6}
\vdeva{राज्यं सुताः कलत्राणि शरीराणि सुखानि च । संसक्तस्यापि नष्टानि तव जन्मनि जन्मनि ॥}
\vtranslit{раджям сутах калатрані шарірані сукхані ча | самсактасяпі наштані тава джанмані джанмані ||}
\vperek{Царство, сини, дружини, тіла й утіхи — усе це вже гинуло для тебе, прив'язаного, — з народження в народження.}
\vcom{Построкове зіставлення: редакторська чернетка, не рецензовано.}
\vcom{Рядок а}
\vdeva{राज्यं सुताः कलत्राणि शरीराणि सुखानि च}
\vtranslit{раджьям сутах калатрані шарірані сукхані ча}
\vperek{Царство, сини, дружини, тіла й утіхи —}
\vcom{rājyaṃ sutāḥ kalatrāṇi śarīrāṇi sukhāni ca перелічує царство, дітей, дружин, тіла й утіхи. Початок із царства враховує становище співрозмовника, але не встановлює зовнішніх історичних дат.}
\vcom{Рядок б}
\vdeva{संसक्तस्यापि नष्टानि तव जन्मनि जन्मनि}
\vtranslit{самсактасьяпі наштані тава джанмані джанмані}
\vperek{усе це вже гинуло для тебе, прив'язаного, — з народження в народження.}
\vcom{saṃsaktasyāpi naṣṭāni tava janmani janmani підкреслює втрати для прив'язаного в повторюваних народженнях. Українське «вже гинуло» передає минулу форму naṣṭāni, не проголошуючи переродження емпірично доведеним фактом.}
\vcom{`naṣṭāni` — доконаний дієприкметник: не «загине», а «гинуло вже не раз». Переклад має зберегти цю повторюваність, інакше вірш звучить як погроза, а він — як нагадування. Звертання адресоване цареві, і царство названо **першим** у переліку: порядок не випадковий.}
\vpar{Українська паралель (ПОТРЕБУЄ ПЕРЕВІРКИ)}{українська поезія XIX ст., не визначено.}{ПОТРЕБУЄ ПЕРЕВІРКИ: думка, що все набуте вже гинуло не раз, трапляється в українській поезії, але називати твір без звірки ми не будемо.}
\vfig{ilustratsii/10-6-motion.pdf}{Низка набутків, що вже гинули з народження в народження}{Схема наша, не історична ілюстрація.}

\vsec{10.7}
\vdeva{अलम् अर्थेन कामेन सुकृतेनापि कर्मणा । एभ्यः संसारकान्तारे न विश्रान्तम् अभून् मनः ॥}
\vtranslit{алам артхена камена сукрітенапі кармана | ебгях самсаракантаре на вішрантам абгун манах ||}
\vperek{Досить уже й багатства, і бажання, і навіть добрих учинків: у цьому лісі сансари твій розум так і не знав відпочинку.}
\vcom{Построкове зіставлення: редакторська чернетка, не рецензовано.}
\vcom{Рядок а}
\vdeva{अलम् अर्थेन कामेन सुकृतेनापि कर्मणा}
\vtranslit{алам артхена камена сукрітенапі кармана}
\vperek{Досить уже й багатства, і бажання, і навіть добрих учинків:}
\vcom{alam arthena kāmena sukṛtenāpi karmaṇā об'єднує багатство, бажання і навіть заслуги/добрі вчинки в одному заклику «досить». sukṛta не слід ототожнювати із запереченням людяності чи допомоги іншим.}
\vcom{Рядок б}
\vdeva{एभ्यः संसारकान्तारे न विश्रान्तम् अभून् मनः}
\vtranslit{ебгьях самсаракантаре на вішрантам абгун манах}
\vperek{у цьому лісі сансари твій розум так і не знав відпочинку.}
\vcom{ebhyaḥ saṃsārakāntāre na viśrāntam abhūn manaḥ містить образ лісу сансари та розуму без відпочинку. Для українського читача ліс — впізнавана образна ситуація, але ми не приписуємо її запозиченню з українського фольклору.}
\vcom{`sukṛta` тут — «заслуга, добре зроблене», а не «добро» взагалі; і текст відкидає саме заслугу, ставлячи її в один ряд із багатством і бажанням. Це найнезручніше місце глави для читача, який звик, що добрі вчинки беззастережно рахуються. `kāntāra` — «ліс», образ блукання, а не місце. `alam` — «годі вже», з інтонацією втоми, а не заборони.}
\vpar{Українська паралель (ПОТРЕБУЄ ПЕРЕВІРКИ)}{українська проза XX ст., не визначено.}{ПОТРЕБУЄ ПЕРЕВІРКИ: образ лісу сансари, у якому розум не знав відпочинку, близький до української прози про виснаження, але конкретного твору ми не перевірили.}
\vfig{ilustratsii/10-7-flame.pdf}{Полумʼя без відпочинку: ліс сансари}{Схема наша, не історична ілюстрація.}

\vsec{10.8}
\vdeva{कृतं न कति जन्मानि कायेन मनसा गिरा । दुःखम् आयासदं कर्म तद् अद्याप्य् उपरम्यताम् ॥}
\vtranslit{крітам на каті джанмані каєна манаса ґіра | духкхам аясадам карма тад адяпй упарам'ятам ||}
\vperek{Скільки ж за ці народження зроблено тілом, розумом і мовою діл, що дають біль і втому! Тож нехай це припиниться хоч тепер.}
\vcom{Построкове зіставлення: редакторська чернетка, не рецензовано.}
\vcom{Рядок а}
\vdeva{कृतं न कति जन्मानि कायेन मनसा गिरा}
\vtranslit{крітам на каті джанмані кайена манаса ґіра}
\vperek{Скільки ж за ці народження зроблено тілом, розумом і мовою діл, що дають біль і втому!}
\vcom{kṛtaṃ na kati janmāni kāyena manasā girā поєднує повторювані народження з трьома інструментальними формами «тілом, розумом, мовою». Український переклад переносить згадку про біль та втому з другого санскритського піврядка в перше речення; сегментація смислова, не побуквена.}
\vcom{Рядок б}
\vdeva{दुःखम् आयासदं कर्म तद् अद्याप्य् उपरम्यताम्}
\vtranslit{духкхам аясадам карма тад адьяпь упарамьятам}
\vperek{Тож нехай це припиниться хоч тепер.}
\vcom{duḥkham āyāsadaṃ karma називає дію, що приносить біль і втому, tad adyāpy uparamyatām подає побажання її припинення. Даний наказ не вимагає стверджувати історичну відмову від праці конкретної особи.}
\vcom{Потрійний інструментал — тіло, розум, мова — це класична формула трьох дверей дії, і переклад має зберегти всі три, а не згортати їх у «вчинки». `uparamyatām` — пасивний імператив: не «припини», а «нехай припиниться»; дія знімається не зусиллям того, хто діє. `adya api` — «хоч тепер», з відтінком докору, а не нейтрального «зараз». Цим віршем глава 10 закінчується.}
\vpar{Українська паралель (наша інтерпретація)}{Василь Стефаник, «Камінний хрест» (1900).}{Наша паралель: у Стефаника втома від зробленого теж стає останнім аргументом людини. Наше порівняння — типологічне, не про джерело.}
\vfig{ilustratsii/10-8-neti.pdf}{Діла тілом, розумом і мовою відкинуті}{Схема наша, не історична ілюстрація.}

